%% ma: L11-B09-0014
%% lop: 11
%% chuong: 3
%% bai: 9
%% dang: 11.9.7
%% loai: TN4
%% muc: NB
%% dap_an: "C"
%% nguon: "de-on-tap-giua-ki-1-toan-11-nam-2026-2027-truong-thpt-ngo-quyen-da-nang.pdf (M09, Đề ôn tập giữa kì 1 Toán 11 THPT Ngô Quyền, Đà Nẵng), Mã đề 0113, Phần I Câu 8"
%% kien_thuc: "Bài 9 (các số đặc trưng đo xu thế trung tâm)"
%% trang_thai: cho-duyet
%% ly_do: "Dạng mới đề xuất 11.9.7"
\begin{ex}
Điểm kiểm tra giữa học kì I của lớp 11A được thống kê theo bảng số liệu ghép nhóm sau:
\begin{center}\begin{tabular}{|l|c|c|c|c|c|}\hline
Số điểm & $[0;2)$ & $[2;4)$ & $[4;6)$ & $[6;8)$ & $[8;10]$\\\hline
Số học sinh & $1$ & $3$ & $8$ & $18$ & $10$\\\hline
\end{tabular}\end{center}
Trong các nhóm sau đây, nhóm nào chứa mốt của bảng số liệu ghép nhóm trên?
\choice{$[8;10]$}{$[4;6)$}{\True $[6;8)$}{$[2;4)$}
\loigiai{Nhóm có tần số lớn nhất ($18$) là $[6;8)$.}
\end{ex}
